% SPDX-License-Identifier: AGPL-3.0-or-later
% Copyright (C) 2026 Sanjeeva Wijeyesakere
%
% XPreview+ User Manual.
% Build:  tectonic xpreview-plus-manual.tex      (or: latexmk -xelatex / -pdf)
\documentclass[11pt,a4paper]{article}

% ---------------------------------------------------------------- engine/fonts
\usepackage{iftex}
\ifXeTeX
  \usepackage{fontspec}
  \setmainfont{texgyrepagella}[Extension=.otf, UprightFont=*-regular, BoldFont=*-bold,
                               ItalicFont=*-italic, BoldItalicFont=*-bolditalic]
  \setsansfont{texgyreheros}[Extension=.otf, UprightFont=*-regular, BoldFont=*-bold,
                             ItalicFont=*-italic, BoldItalicFont=*-bolditalic, Scale=0.92]
  \setmonofont{texgyrecursor}[Extension=.otf, UprightFont=*-regular, BoldFont=*-bold, Scale=0.95]
\else
  \usepackage[T1]{fontenc}
  \usepackage[utf8]{inputenc}
  \usepackage{tgpagella}
  \usepackage[scale=0.92]{tgheros}
  \usepackage{tgcursor}
\fi

% ---------------------------------------------------------------- packages
\usepackage[margin=2.4cm,headheight=14pt]{geometry}
\usepackage{graphicx}
\usepackage[table]{xcolor}
\usepackage{booktabs,tabularx,longtable,array}
\usepackage{enumitem}
\usepackage{microtype}
\usepackage{fancyhdr}
\usepackage[most]{tcolorbox}
\usepackage{titlesec}
\usepackage[section]{placeins}  % keep figures inside their section
\usepackage[hidelinks]{hyperref}
\usepackage{bookmark}

\definecolor{brand}{HTML}{1F5FBF}
\definecolor{brandlight}{HTML}{E8F0FB}
\definecolor{notebg}{HTML}{FFF6D6}
\definecolor{noteframe}{HTML}{E2C766}
\definecolor{keybg}{HTML}{F4F4F6}
\definecolor{keyframe}{HTML}{B9BBC2}
\definecolor{rule}{HTML}{C9CED8}

\hypersetup{colorlinks=true, linkcolor=brand, urlcolor=brand,
  pdftitle={XPreview+ User Manual}, pdfauthor={Sanjeeva Wijeyesakere},
  pdfsubject={User manual for XPreview+ 2.1}}

\setlist{itemsep=2pt, topsep=4pt}
\setlength{\parindent}{0pt}
\setlength{\parskip}{6pt}
\renewcommand{\arraystretch}{1.25}

\titleformat{\section}{\Large\sffamily\bfseries\color{brand}}{\thesection}{0.8em}{}
\titleformat{\subsection}{\large\sffamily\bfseries}{\thesubsection}{0.8em}{}
\titleformat{\subsubsection}{\normalsize\sffamily\bfseries}{}{0em}{}
\titlespacing*{\section}{0pt}{18pt}{8pt}
\titlespacing*{\subsection}{0pt}{12pt}{4pt}
\titlespacing*{\subsubsection}{0pt}{10pt}{2pt}
\setcounter{tocdepth}{2}
\usepackage{tocloft}  % compact table of contents (fits one page)
\setlength{\cftbeforesecskip}{3.5pt}
\renewcommand{\cftsecfont}{\sffamily\bfseries}
\renewcommand{\cftsecpagefont}{\sffamily\bfseries}
\setlength{\cftbeforetoctitleskip}{0pt}
\renewcommand{\cfttoctitlefont}{\Large\sffamily\bfseries\color{brand}}

\pagestyle{fancy}
\fancyhf{}
\fancyhead[L]{\small\sffamily XPreview+ 2.1 User Manual}
\fancyhead[R]{\small\sffamily\nouppercase{\leftmark}}
\fancyfoot[C]{\small\sffamily\thepage}
\renewcommand{\headrulewidth}{0.4pt}
\renewcommand{\sectionmark}[1]{\markboth{#1}{}}

% ---------------------------------------------------------------- macros
\newcommand{\app}{XPreview+}
% Keyboard key cap, e.g. \key{Cmd}+\key{O}
\newcommand{\key}[1]{{\setlength{\fboxsep}{1.4pt}\setlength{\fboxrule}{0.5pt}%
  \raisebox{-0.1ex}{\fcolorbox{keyframe}{keybg}{\footnotesize\sffamily\strut #1}}}}
% Menu path, e.g. \menu{File > Export as PDF...}
\newcommand{\menu}[1]{\textsf{\textbf{#1}}}
\newcommand{\ui}[1]{\textsf{#1}}
\newcommand{\file}[1]{\texttt{#1}}

\newtcolorbox{note}[1][Note]{enhanced, breakable, colback=notebg, colframe=noteframe,
  boxrule=0.5pt, arc=2pt, left=6pt, right=6pt, top=4pt, bottom=4pt,
  fonttitle=\sffamily\bfseries\small, coltitle=black, attach title to upper={\ }, title={#1:}}
\newtcolorbox{tip}{enhanced, breakable, colback=brandlight, colframe=brand!45, boxrule=0.5pt,
  arc=2pt, left=6pt, right=6pt, top=4pt, bottom=4pt,
  fonttitle=\sffamily\bfseries\small, coltitle=black, attach title to upper={\ }, title={Tip:}}

\newcommand{\screenshot}[3]{% file, width, caption
  \begin{figure}[htbp]\centering
    % The frame adds 2*(\fboxsep+\fboxrule); keep the total within #2.
    \fcolorbox{rule}{white}{\includegraphics[width=\dimexpr#2-2\fboxsep-2\fboxrule\relax]{figures/#1}}
    \caption{#3}\label{fig:#1}
  \end{figure}}

\begin{document}

% ================================================================ title page
\begin{titlepage}
  \centering
  \vspace*{2.5cm}
  \includegraphics[width=4.2cm]{figures/icon.png}\par
  \vspace{1.2cm}
  {\sffamily\bfseries\fontsize{38}{44}\selectfont \app\par}
  \vspace{0.4cm}
  {\sffamily\Large User Manual\par}
  \vspace{0.3cm}
  {\sffamily\large Version 2.1\par}
  \vspace{1.6cm}
  {\large A fast viewer for PDF documents, images, RTF and OpenDocument text,\\
  with PDF form filling, page editing, and export of anything to PDF.\par}
  \vfill
  {\sffamily macOS \quad·\quad Windows \quad·\quad Linux\par}
  \vspace{0.6cm}
  {\small Copyright \textcopyright{} 2026 Sanjeeva Wijeyesakere.\\
  \app{} is free software under the GNU Affero General Public License, version 3 or later.\\
  \url{https://github.com/swijeyes/XPreview-Plus}\par}
\end{titlepage}

\tableofcontents
\clearpage

% ================================================================
\section{Introduction}

\app{} is a document and image viewer in the spirit of a built-in document viewer. It opens a
file almost instantly, shows it page by page, lets you search, select and copy
text, and can turn any file it opens into a PDF. PDFs can also be edited: fill in their forms,
and rearrange, rotate, delete and combine their pages. It runs on macOS and Windows, and can be built for Linux.

\subsection{What \app{} can open}
\begin{itemize}
  \item \textbf{PDF documents}, including password-protected ones.
  \item \textbf{Images}: PNG, JPEG, GIF, WebP, TIFF (including multi-page TIFF), BMP,
        SVG and more; HEIC on macOS. See Section~\ref{sec:formats} for the complete list.
  \item \textbf{Rich Text Format} (\file{.rtf}) documents, such as those from TextEdit, WordPad or Word.
  \item \textbf{OpenDocument Text} (\file{.odt}, \file{.ott}, \file{.fodt}) documents,
        such as those from LibreOffice, Google Docs or Microsoft Word.
\end{itemize}

\subsection{Highlights}
\begin{itemize}
  \item \textbf{Page thumbnails} and a \textbf{table of contents} in the sidebar.
  \item \textbf{Export as PDF} for every file type.
  \item \textbf{Fill in PDF forms}: text fields, check boxes, radio buttons and lists.
  \item \textbf{Edit PDF pages}: drag thumbnails to reorder pages, drop other PDFs (or images and
        documents) into the sidebar to combine files, and rotate, delete or insert pages.
  \item \textbf{Search} with highlighted results, \textbf{text selection} and \textbf{copy}.
  \item Smooth \textbf{zooming} (including trackpad pinch), \textbf{fit to width or page}, and \textbf{rotation}.
  \item \textbf{Tabs}, so several documents can be open in one window.
  \item \textbf{Private and safe}: \app{} never connects to the internet on its own, never runs
        macros or scripts, and treats every file as potentially hostile.
\end{itemize}

\screenshot{window-odt.png}{\linewidth}{\app{} showing an OpenDocument text file, with page
thumbnails in the sidebar.}

% ================================================================
\section{Installing \app{}}

\subsection{macOS}
\textbf{Requirements:} a Mac with Apple silicon (M-series) and a recent version of macOS.

\begin{enumerate}
  \item Download \file{XPreviewPlus-2.1-macos-arm64.dmg} from the project's
        \href{https://github.com/swijeyes/XPreview-Plus/releases/latest}{Releases page}.
  \item Double-click the downloaded file to open it.
  \item Drag the \ui{\app} icon onto the \ui{Applications} folder shown next to it.
  \item Eject the disk image (drag it to the Trash, or click the eject button in Finder).
\end{enumerate}

\subsubsection{Opening it the first time}
If the copy you downloaded is not notarised by Apple, macOS blocks it the first time
you open it and shows a message that it ``cannot verify'' the app. To allow it:
\begin{enumerate}
  \item Try to open \app{} once and close the warning.
  \item Open \menu{System Settings > Privacy \& Security}.
  \item Scroll down to the message about \app{} and click \ui{Open Anyway}.
  \item Confirm with your password or Touch ID.
\end{enumerate}
You only need to do this once.

\subsubsection{Removing \app{}}
Drag \ui{\app} from \ui{Applications} to the Trash. To also remove its settings,
see Section~\ref{sec:settings}.

\subsection{Windows}
\textbf{Requirements:} 64-bit Windows 10 or Windows 11.

\app{} for Windows is \emph{portable}. There is no installer, and it never writes to the
Windows registry.
\begin{enumerate}
  \item Download \file{XPreviewPlus-2.1-windows-x64-portable.zip} from the
        \href{https://github.com/swijeyes/XPreview-Plus/releases/latest}{Releases page}.
  \item Right-click the zip file, choose \ui{Extract All\ldots}, and pick a folder,
        for example your \file{Documents} folder or a USB stick.
  \item Open the extracted \file{XPreview+} folder and double-click \file{XPreview+.exe}.
\end{enumerate}

If Microsoft Defender SmartScreen shows ``Windows protected your PC'', click
\ui{More info} and then \ui{Run anyway}. This appears for new applications that are not yet
widely downloaded.

\begin{note}
Keep all the files in the \file{XPreview+} folder together; the program needs them.
Your settings are stored in a \file{settings} folder inside it. To remove \app{}, delete the folder.
\end{note}

\subsection{Linux}
There is no ready-made Linux package yet. \app{} can be built from source; see
the \file{README} in the project repository.

\subsection{Checking your download (optional)}
Each download on the Releases page has a matching \file{.sha256} file containing its
checksum. To check a download, open a terminal in the folder where you saved both files and run
\texttt{shasum -a 256 -c} followed by the name of the \file{.sha256} file (on Windows,
\texttt{CertUtil -hashfile \textit{file} SHA256}), and compare.

% ================================================================
\section{The \app{} window}\label{sec:window}

\screenshot{window-pdf.png}{\linewidth}{The main window with a PDF open. Thumbnails are on
the left, the page view on the right, and the status bar along the bottom.}

\subsection{Toolbar}
From left to right:
\begin{description}[style=nextline, leftmargin=2.2em, font=\sffamily\bfseries]
  \item[Sidebar button] Click to show or hide the sidebar. Click the small arrow next to it
        to choose between \ui{Thumbnails} and \ui{Table of Contents}.
  \item[Zoom out / Zoom in] Step the zoom down or up.
  \item[Rotate left / Rotate right] Turn the pages a quarter turn.
  \item[Export as PDF] Save the current document as a PDF (Section~\ref{sec:export}).
  \item[Search field] Find text in the document (Section~\ref{sec:search}). The number of
        results appears to its left.
\end{description}

\subsection{Sidebar}
The sidebar shows either \textbf{page thumbnails} or the \textbf{table of contents}. Drag the
thin divider between the sidebar and the page to make it wider or narrower; thumbnails grow
and shrink with it. For PDFs, the thumbnails are also where you rearrange and combine pages
(Section~\ref{sec:pages}).

\subsection{Page view}
The main area shows the pages of your document one after another on a grey background. Scroll to move
through the document.

\subsection{Status bar}
The left side shows the current page, for example ``Page 2 of 9''. For single-page files it shows the file type instead.
The right side has the \textbf{zoom box}. It shows the current zoom, and you can
choose \ui{Fit Width}, \ui{Fit Page} or a percentage from it, or type a value such as
\texttt{140\%} and press \key{Return}.

\subsection{Information bar}
For RTF and ODT documents that \app{} had to simplify, a yellow bar appears above the page
(Section~\ref{sec:simplified}). For PDFs with fillable form fields, the bar says so
(Section~\ref{sec:forms}). Click the \ui{×} button on its right to dismiss it.

\subsection{Tabs}
Each file you open gets its own tab. The tab bar appears as soon as two or more files are open.
Click a tab to switch, drag tabs to reorder them, and click a tab's close button
(or press \key{Cmd}+\key{W}) to close it.

% ================================================================
\section{Opening files}

There are several ways to open a file:
\begin{itemize}
  \item Choose \menu{File > Open\ldots} (\key{Cmd}+\key{O} on macOS, \key{Ctrl}+\key{O} on
        Windows) and select one or more files.
  \item \textbf{Drag} files onto the \app{} window.
  \item On macOS, drag files onto the \app{} icon in the \textbf{Dock}.
  \item In Finder, right-click a file and choose \menu{Open With > \app}.
  \item Choose a recently opened file from \menu{File > Open Recent}. The list keeps the last
        twelve files; \ui{Clear Menu} empties it.
\end{itemize}

If a file is already open, \app{} switches to its tab instead of opening it twice.

\begin{tip}
To have a type of file always open in \app{} on macOS, select such a file in Finder, choose
\menu{File > Get Info}, pick \app{} under \ui{Open with}, and click \ui{Change All\ldots}.
\end{tip}

\subsection{Password-protected PDFs}
When a PDF is encrypted, \app{} asks for its password. All standard PDF encryption is
supported, from older RC4 files to AES-256. If the password is wrong, \app{} says so and asks
again; each further try waits a little longer. If you cancel, the tab shows a message instead
of the document; close it and open the file again to retry.

Some PDFs open without a password but carry \textbf{restrictions} set by their owner: no
printing, no copying, or no changes. \app{} honours them --- the yellow bar says what is
restricted, and the matching commands are unavailable. Opening the file with its
\emph{owner} password lifts the restrictions. Section~\ref{sec:protect} explains how to protect
your own PDFs.

\subsection{PostScript and EPS files}
\app{} opens PostScript (\file{.ps}) and Encapsulated PostScript (\file{.eps}) files by turning
them into a PDF with Ghostscript, which is included in the macOS app. EPS images are cropped to
their bounding box. The converted document can be searched, copied and printed like a PDF;
to keep it as a PDF, use \menu{File > Save As\ldots}. The original file is never changed.

\begin{note}
PostScript files are programs. Ghostscript runs them in a separate, locked-down process: it
cannot use the network, start other programs, read your documents, or write anywhere except
its own temporary folder, and it is stopped after a minute. On Windows, PostScript and EPS
open if Ghostscript is installed.
\end{note}

\subsection{How \app{} recognises files}
\app{} looks at a file's \emph{contents}, not just its name, to decide what it is. If a file's
extension is wrong (for example an RTF document named \file{report.pdf}), \app{} still opens it
correctly and notes ``detected as \ldots'' in the status bar.

\subsection{Files that cannot be opened}
If a file is damaged, of an unsupported type, or too large to open safely, its tab shows a short explanation.
OpenDocument spreadsheets (\file{.ods}) and presentations
(\file{.odp}) are recognised but not supported in this version.

% ================================================================
\section{Viewing and navigating}

\subsection{Moving through pages}
\begin{itemize}
  \item Scroll with the trackpad, mouse wheel or scroll bar.
  \item Click a \textbf{thumbnail} to jump to that page. The highlighted thumbnail follows
        the page you are reading.
  \item Click an entry in the \textbf{table of contents}.
  \item Use the \menu{Go} menu: \ui{Previous Page}, \ui{Next Page}, \ui{First Page},
        \ui{Last Page} and \ui{Go to Page\ldots}.
  \item On the keyboard: \key{Space} and \key{Shift}+\key{Space} move a screen down or up;
        \key{Page Up}/\key{Page Down} and the arrow keys scroll; \key{Home} and \key{End} go to the
        first and last page. When the page fits the window's width, \key{←} and \key{→} move to
        the previous and next page.
\end{itemize}

\subsection{Table of contents}
For PDFs, the table of contents comes from the document's bookmarks. For RTF and ODT
documents, \app{} builds it from the headings. If a document has neither, the sidebar shows
``No table of contents''.

\subsection{Zooming}
\begin{itemize}
  \item Use the zoom buttons in the toolbar, or \menu{View > Zoom In} / \menu{Zoom Out}.
  \item \menu{View > Actual Size} shows the page at 100\%.
  \item \menu{View > Zoom to Fit} fits the whole page in the window; \menu{View > Fit Width}
        fits the page width. Both keep fitting as you resize the window.
  \item On a trackpad, \textbf{pinch} to zoom. Double-tap with two fingers to switch between fit
        width and a closer view.
  \item Hold \key{Cmd} (macOS) or \key{Ctrl} (Windows) and scroll to zoom around the pointer.
  \item Use the zoom box in the status bar.
\end{itemize}
Zoom ranges from 5\% to 3200\%. Documents open fitted to the window's width. Images open
fitted to the window, but are never enlarged beyond 100\%.

\subsection{Rotating}
\menu{Tools > Rotate Left} and \menu{Rotate Right} (\key{Cmd}+\key{L}, \key{Cmd}+\key{R}) turn
all pages. Rotation only affects how the document is shown; the original file is never changed.
Exported PDFs and printouts use the rotation you see.

\subsection{Viewing images}
\screenshot{window-image.png}{0.8\linewidth}{An SVG image. Images are a single page and are
shown at their real size if they fit.}

When an image is larger than the window, \textbf{drag} it to move it around. \textbf{Double-click}
to switch between fitting the window and actual size. Multi-page TIFF files show
one page per image. Animated GIF and WebP images show their first frame.

% ================================================================
\section{Searching}\label{sec:search}

\begin{enumerate}
  \item Click the search field in the toolbar, or press \key{Cmd}+\key{F} (\key{Ctrl}+\key{F}).
  \item Type the text you are looking for and press \key{Return}.
\end{enumerate}
Every match is highlighted in yellow and the current one in orange. The counter next to the
field shows, for example, ``3 of 12'', or ``Not found''.

\begin{itemize}
  \item Press \key{Return} again, or choose \menu{Edit > Find Next}, to go to the next match.
  \item \key{Shift}+\key{Return} or \menu{Edit > Find Previous} goes back.
  \item Searching ignores upper and lower case.
  \item Clear the field (click its \ui{×}) to remove the highlights.
\end{itemize}

\begin{note}
Images have no text to search. Scanned PDFs that contain only pictures of pages cannot be
searched either, because \app{} does not perform text recognition (OCR).
\end{note}

% ================================================================
\section{Selecting and copying}

\subsection{Text}
In PDF, RTF and ODT documents, drag across text to select it; the selection is shown in blue.
Then choose \menu{Edit > Copy} (\key{Cmd}+\key{C}) to copy it. A selection covers text on one
page at a time. Click anywhere or press \key{Esc} to clear it.

\begin{tip}
To move around a zoomed-in document instead of selecting text, hold \key{Option} (\key{Alt} on
Windows) while you drag.
\end{tip}

\subsection{Images}
When an image is open and nothing is selected, \menu{Edit > Copy} copies the image itself, ready to paste
into another application.

\subsection{Links}
Click a link to follow it. Web and e-mail links (\texttt{http}, \texttt{https} and
\texttt{mailto}) open in your browser or mail app. Links inside a PDF jump to the page they
point to. The pointer turns into a hand over a link. For safety, other kinds of links are
ignored, and nothing is ever opened without a click.

% ================================================================
\section{Exporting as PDF}\label{sec:export}

Any file that \app{} can open can be saved as a PDF.
\begin{enumerate}
  \item Choose \menu{File > Export as PDF\ldots} (\key{Shift}+\key{Cmd}+\key{E}), or click the
        export button in the toolbar.
  \item Choose a name and location. \app{} suggests the original name with \file{.pdf}, in the same folder.
  \item Click \ui{Save}.
\end{enumerate}
A progress window appears for longer documents, and the status bar confirms when the export
is done.

\subsection{What you get}
\begin{description}[style=nextline, leftmargin=2.2em, font=\sffamily\bfseries]
  \item[PDF documents] An exact copy of the original file. If you rotated the view,
        the copy is saved with that rotation; if you filled in fields or edited pages, the
        copy includes those changes.
  \item[RTF and ODT documents] Real PDF pages with sharp, selectable and searchable text, the
        same layout you see in \app{}, and the fonts embedded.
  \item[Images] A PDF page the size of the image, with the picture at full quality. SVG images
        stay sharp at any zoom.
\end{description}

\begin{note}
Exports are written safely: the new PDF only replaces an existing file once it has been
written completely, so a failed or interrupted export never leaves a broken file behind.
\end{note}

% ================================================================
\section{Filling in PDF forms}\label{sec:forms}

Many PDFs --- applications, tax forms, registration sheets --- contain \textbf{form fields}.
\app{} shows these fields with a light blue tint, and a yellow bar above the page tells you
that the document has fields to fill in.

\begin{description}[style=nextline, leftmargin=2.2em, font=\sffamily\bfseries]
  \item[Text fields] Click the field and type. The pointer becomes a text cursor over a text
        field. Use the arrow keys, \key{Backspace} and \key{Delete} as usual; drag or
        double-click to select text. \key{Cmd}+\key{C}, \key{Cmd}+\key{X}, \key{Cmd}+\key{V},
        \key{Cmd}+\key{A} and \key{Cmd}+\key{Z} work inside the field.
  \item[Check boxes and radio buttons] Click to tick or untick. With the field selected,
        \key{Space} toggles it too.
  \item[Lists and pop-up menus] Click to open the list, then click a choice.
\end{description}

Press \key{Tab} to move to the next field on the page, and \key{Shift}+\key{Tab} to go back.
Press \key{Esc}, or click anywhere outside the fields, to finish editing a field.

\screenshot{window-form.png}{\linewidth}{A PDF form being filled in. Fields have a light
blue tint.}

When you have filled in the form, save it (Section~\ref{sec:saving}). The values are stored in
the PDF itself, so any other PDF application shows them.

\begin{note}
\app{} never runs scripts in forms, so fields that calculate totals or check what you type in
other applications do not do so in \app{}. Digital signature fields cannot be signed.
\end{note}

% ================================================================
\section{Editing PDF pages}\label{sec:pages}

The page thumbnails of a PDF are its pages: what you change there changes the document.
Click a thumbnail to select a page; \key{Shift}-click to select a range, or \key{Cmd}-click
(\key{Ctrl}-click on Windows and Linux) to add single pages. The commands below act on the
selected pages --- or on the current page if no thumbnail is selected.

\subsection{Rearranging pages}
Drag one or more selected thumbnails up or down. A blue line shows where they will go; release
the mouse button to move them there. Hold \key{Option} (\key{Alt}) while you drop to put a
\emph{copy} of the pages there instead.

\subsection{Combining PDFs}
Drag a PDF from the Finder (or Explorer) onto the thumbnails. The blue line shows where its
pages will be inserted; drop it, and all of its pages are added there. You can drop several
files at once. Images, PostScript, EPS, RTF and ODT documents can be dropped too: they are turned
into PDF pages first, exactly as \menu{File > Export as PDF} would.

You can also drag thumbnails from one \app{} window into another to copy pages between
documents, or choose \menu{Edit > Insert > Page from File\ldots} to pick files from a dialog.
The new pages go after the selected page.

\begin{tip}
Dropping a file on the \emph{page view} (rather than the thumbnails) opens it in its own tab
as usual.
\end{tip}

\subsection{Rotating, deleting and adding pages}
\begin{description}[style=nextline, leftmargin=2.2em, font=\sffamily\bfseries]
  \item[Rotate] \menu{Tools > Rotate Page Left} or \menu{Rotate Page Right}
        (\key{Shift}+\key{Cmd}+\key{L} / \key{Shift}+\key{Cmd}+\key{R}) turns the selected
        pages and is saved with the document. (\menu{Rotate Left} and \menu{Rotate Right},
        without ``Page'', only turn your view of the whole document.)
  \item[Delete] Press \key{Delete} or \key{Backspace} in the thumbnails, or choose
        \menu{Edit > Delete Page}. A PDF must keep at least one page.
  \item[Blank page] \menu{Edit > Insert > Blank Page} adds an empty page, the same size as the
        current one, after it.
  \item[Export pages] \menu{Tools > Export Selected Pages\ldots} saves just the selected pages
        as a new PDF.
\end{description}
All of these are also in the menu that appears when you right-click (or \key{Ctrl}-click) a
thumbnail.

\subsection{Undo}
\menu{Edit > Undo} (\key{Cmd}+\key{Z}) takes back the last page change, and
\menu{Edit > Redo} (\key{Shift}+\key{Cmd}+\key{Z}) puts it back. The menu names the change,
for example \menu{Undo Move Page}. The last 20 page changes can be undone.

\begin{note}
When you combine PDFs, the pages you insert keep their look, but form fields on them can no
longer be filled in. Fill in a form before combining it with other files.
\end{note}

% ================================================================
\section{Saving your changes}\label{sec:saving}

A PDF you have changed is marked as edited: on macOS the window title ends in
``--- Edited'' and the close button shows a dot; on Windows and Linux the title has an
asterisk. The tab of an edited document shows a dot (•) after its name.
\begin{description}[style=nextline, leftmargin=2.2em, font=\sffamily\bfseries]
  \item[Save] \menu{File > Save} (\key{Cmd}+\key{S}) writes the changes into the PDF.
  \item[Save As] \menu{File > Save As\ldots} (\key{Shift}+\key{Cmd}+\key{S}) writes them to a
        new file and leaves the original as it was. The window then shows the new file.
  \item[Export as PDF] also includes your changes, without changing which file is open.
\end{description}
If you close a tab or window, or quit, with unsaved changes, \app{} asks whether to save them.
Like exports, saving is safe: the file is only replaced once the new version has been written
completely.

% ================================================================
\section{Password protection}\label{sec:protect}

Choose \menu{File > Password Protection\ldots} to protect a PDF. You can:
\begin{description}[style=nextline, leftmargin=2.2em, font=\sffamily\bfseries]
  \item[Require a password to open it] Nobody can see the document without the password.
  \item[Restrict printing, copying and changes] Others can open the PDF but not print it, copy
        its text, or change its pages and forms. An \textbf{owner password}, which must differ
        from the open password, lifts the restrictions.
\end{description}
The protection uses AES-256, the strongest encryption PDF offers, and takes effect when you
save (\menu{File > Save} or \menu{Save As\ldots}). \menu{Export as PDF} applies it too.

When you edit and save a PDF that was already protected, \app{} keeps its protection as it
was. To change or remove it, open the PDF with its owner password (or its only password),
choose \menu{File > Password Protection\ldots}, then \ui{OK} or \ui{Remove Protection}, and save.

\begin{note}
Keep your password safe: a PDF protected with a forgotten password cannot be opened, by
\app{} or anyone else. Restrictions are honoured by PDF applications, but are not a
substitute for an open password.
\end{note}

% ================================================================
\section{Printing}

Choose \menu{File > Print\ldots} (\key{Cmd}+\key{P}). In the print dialog you can pick a printer
and the range of pages. Each page is scaled to fit the paper and centred, using the current
rotation.

% ================================================================
\section{Document information and the ``simplified view''}\label{sec:simplified}

\subsection{The Inspector}
\menu{Tools > Show Inspector} (\key{Cmd}+\key{I}) shows information about the open file: its type,
the number of pages and details such as title, author and dates. For images it shows the
dimensions and colour profile. It also lists anything that was shown in a simplified way.

\subsection{Simplified view}
RTF and ODT files are word-processor documents. \app{} contains its own renderer for them, so
you do not need LibreOffice or Microsoft Word installed. It reproduces text, fonts, colours, lists,
tables, images, headers and footers faithfully, but a few features are approximated. When
that happens, a yellow bar says
\emph{``Simplified view: some formatting may differ from the original.''} Its buttons are:
\begin{description}[style=nextline, leftmargin=2.2em, font=\sffamily\bfseries]
  \item[Details\ldots] Lists exactly what was simplified, for example ``Footnotes are shown at
        the end of the document'' or ``Font substituted: Calibri → Carlito''.
  \item[Open in Default App] Opens the file in the application your computer normally uses for it.
  \item[Export as PDF\ldots] Saves the document as you see it.
\end{description}
\menu{File > Open in Default App} is also available for any file, except one whose name does not
match its contents (for example a PDF named \file{.hta} or \file{.command}): your computer picks the
application by the name, so such a file could run as a program instead of opening as a document.

\subsection{Fonts}
Documents often use fonts that are not installed on your computer. \app{} includes free fonts
with the same letter widths as the common Microsoft fonts, so lines and pages break in the same
places as in the original:

\begin{center}
\begin{tabular}{ll}
\toprule
\textbf{Font in the document} & \textbf{Shown with} \\
\midrule
Arial, Helvetica & Arimo \\
Times New Roman, Times & Tinos \\
Courier New, Courier & Cousine \\
Calibri & Carlito \\
Cambria & Caladea \\
\bottomrule
\end{tabular}
\end{center}
A font that is installed on your computer is always used as is.

% ================================================================
\section{Supported file types}\label{sec:formats}

\begin{longtable}{>{\raggedright\arraybackslash}p{3.3cm}>{\raggedright\arraybackslash}p{3.2cm}>{\raggedright\arraybackslash}p{8.2cm}}
\toprule
\textbf{Type} & \textbf{Extensions} & \textbf{Notes} \\
\midrule
\endhead
PDF & \file{.pdf} & Full support, including encrypted files (RC4 and AES), forms, bookmarks and links. \\
PostScript, EPS & \file{.ps .eps .epsf .epsi} & Converted to PDF with the included Ghostscript (macOS); on Windows if Ghostscript is installed. \\
PNG, JPEG, BMP & \file{.png .jpg .jpeg .bmp} & Rotation stored in photos (EXIF) is applied automatically. \\
GIF, WebP & \file{.gif .webp} & Animated images show their first frame. \\
TIFF & \file{.tif .tiff} & Multi-page TIFF files show all pages. \\
SVG & \file{.svg .svgz} & Vector images; stay sharp at every zoom. \\
HEIC, JPEG 2000 & \file{.heic .heif .jp2} & macOS only. \\
Other images & \file{.ico .icns .tga .wbmp} & \\
Rich Text Format & \file{.rtf} & Fonts, colours, lists, tables, pictures (PNG, JPEG), links, headers, footers, footnotes. \\
OpenDocument Text & \file{.odt .ott .fodt} & Styles, headings, lists, tables, images, links, headers, footers, page numbers, footnotes. \\
OpenDocument Spreadsheet / Presentation & \file{.ods .odp} & Recognised, but not supported yet. \\
\bottomrule
\end{longtable}

\subsection{Known limitations of the RTF and ODT renderer}
\begin{itemize}
  \item Multiple text columns are shown as a single column.
  \item Footnotes are collected at the end of the document rather than at the foot of each page.
  \item Text boxes and frames are placed in the flow of the text, near where they are anchored.
  \item Drawings made of shapes, embedded charts without a preview picture, and WMF/EMF pictures
        (common in older RTF files) appear as labelled placeholders.
  \item All pages use the page size of the document's first section.
  \item Comments are hidden. Documents with tracked changes show the final text.
\end{itemize}

% ================================================================
\section{Keyboard shortcuts}\label{sec:shortcuts}

On macOS, \key{Cmd} is the Command key and \key{Option} is the Option (Alt) key.

\begin{longtable}{>{\raggedright\arraybackslash}p{5.6cm}>{\raggedright\arraybackslash}p{4.6cm}>{\raggedright\arraybackslash}p{4.6cm}}
\toprule
\textbf{Action} & \textbf{macOS} & \textbf{Windows / Linux} \\
\midrule
\endhead
\multicolumn{3}{l}{\sffamily\bfseries File} \\
Open & \key{Cmd}+\key{O} & \key{Ctrl}+\key{O} \\
Close tab & \key{Cmd}+\key{W} & \key{Ctrl}+\key{W} \\
Save (PDF) & \key{Cmd}+\key{S} & \key{Ctrl}+\key{S} \\
Save As (PDF) & \key{Shift}+\key{Cmd}+\key{S} & \key{Ctrl}+\key{Shift}+\key{S} \\
Export as PDF & \key{Shift}+\key{Cmd}+\key{E} & \key{Ctrl}+\key{Shift}+\key{E} \\
Print & \key{Cmd}+\key{P} & \key{Ctrl}+\key{P} \\
Quit & \key{Cmd}+\key{Q} & \key{Ctrl}+\key{Q} on Linux; close the window on Windows \\
\midrule
\multicolumn{3}{l}{\sffamily\bfseries Edit} \\
Undo / redo & \key{Cmd}+\key{Z} / \key{Shift}+\key{Cmd}+\key{Z} & \key{Ctrl}+\key{Z} / \key{Ctrl}+\key{Y} \\
Copy & \key{Cmd}+\key{C} & \key{Ctrl}+\key{C} \\
Insert page from file & \key{Option}+\key{Cmd}+\key{I} & \key{Ctrl}+\key{Alt}+\key{I} \\
Insert blank page & \key{Shift}+\key{Cmd}+\key{N} & \key{Ctrl}+\key{Shift}+\key{N} \\
Delete selected pages & \key{Delete} in the thumbnails & \key{Delete} in the thumbnails \\
Find & \key{Cmd}+\key{F} & \key{Ctrl}+\key{F} \\
Find next & \key{Cmd}+\key{G} & \key{F3} or \key{Ctrl}+\key{G} \\
Find previous & \key{Shift}+\key{Cmd}+\key{G} & \key{Shift}+\key{F3} or \key{Ctrl}+\key{Shift}+\key{G} \\
\midrule
\multicolumn{3}{l}{\sffamily\bfseries View} \\
Hide sidebar & \key{Option}+\key{Cmd}+\key{1} & \key{Ctrl}+\key{Alt}+\key{1} \\
Thumbnails & \key{Option}+\key{Cmd}+\key{2} & \key{Ctrl}+\key{Alt}+\key{2} \\
Table of contents & \key{Option}+\key{Cmd}+\key{3} & \key{Ctrl}+\key{Alt}+\key{3} \\
Actual size & \key{Cmd}+\key{0} & \key{Ctrl}+\key{0} \\
Zoom in & \key{Cmd}+\key{+} or \key{Cmd}+\key{=} & \key{Ctrl}+\key{+} \\
Zoom out & \key{Cmd}+\key{-} & \key{Ctrl}+\key{-} \\
Zoom to fit & \key{Cmd}+\key{9} & \key{Ctrl}+\key{9} \\
Fit width & \key{Cmd}+\key{8} & \key{Ctrl}+\key{8} \\
\midrule
\multicolumn{3}{l}{\sffamily\bfseries Go} \\
Previous / next page & \key{Option}+\key{Cmd}+\key{↑} / \key{↓} & \key{Ctrl}+\key{Alt}+\key{↑} / \key{↓} \\
First / last page & \key{Cmd}+\key{↑} / \key{Cmd}+\key{↓} & \key{Ctrl}+\key{↑} / \key{Ctrl}+\key{↓} \\
Go to page & \key{Option}+\key{Cmd}+\key{G} & \key{Ctrl}+\key{Alt}+\key{G} \\
\midrule
\multicolumn{3}{l}{\sffamily\bfseries Tools} \\
Rotate left / right & \key{Cmd}+\key{L} / \key{Cmd}+\key{R} & \key{Ctrl}+\key{L} / \key{Ctrl}+\key{R} \\
Rotate page left / right & \key{Shift}+\key{Cmd}+\key{L} / \key{R} & \key{Ctrl}+\key{Shift}+\key{L} / \key{R} \\
Show Inspector & \key{Cmd}+\key{I} & \key{Ctrl}+\key{I} \\
\midrule
\multicolumn{3}{l}{\sffamily\bfseries In the page view} \\
Screen down / up & \key{Space} / \key{Shift}+\key{Space} & \key{Space} / \key{Shift}+\key{Space} \\
First / last page & \key{Home} / \key{End} & \key{Home} / \key{End} \\
Previous / next page (page fits width) & \key{←} / \key{→} & \key{←} / \key{→} \\
Clear selection & \key{Esc} & \key{Esc} \\
Zoom around pointer & \key{Cmd}+scroll, or pinch & \key{Ctrl}+scroll \\
Move a zoomed document & \key{Option}+drag & \key{Alt}+drag \\
\bottomrule
\end{longtable}

% ================================================================
\section{Settings and privacy}\label{sec:settings}

\subsection{Preferences}
Open \app{}'s preferences with \key{Cmd}+\key{,} on macOS, or from \menu{Edit > Preferences} on
Windows and Linux. Changes take effect immediately. The window has three tabs:
\begin{description}[style=nextline, leftmargin=2.2em, font=\sffamily\bfseries]
  \item[General] the window's look --- a \textbf{Modern} flat theme in Light, Dark or Automatic
        (Light is the default), or a \textbf{Classic} glossy look in Blue or Gray. For Modern,
        \ui{Dark style} chooses how it looks when dark: neutral \textbf{Graphite}, or
        \textbf{Dracula}, with the purple-grey colours of the popular Dracula theme
        (Figure~\ref{fig:window-dracula.png}); the scroll-bar
        style when using the Modern theme; and whether opening several files keeps them together in
        one window or gives each file its own window.
  \item[PDF] what Actual Size (100\%) shows --- pages at their printed size, or one PDF point per
        screen pixel --- whether to return to the page you were last reading, and the initial page
        layout.
  \item[Images] what Actual Size (100\%) shows for images --- their printed size, using the
        resolution saved in the image, or one image pixel per screen pixel.
\end{description}

\subsection{What \app{} remembers}
Separately from the settings above, \app{} remembers:
\begin{itemize}
  \item the list of recently opened files,
  \item the sidebar mode (hidden, thumbnails or table of contents),
  \item the window's size and position (saved as soon as you move or resize it, and
        used for the next window you open),
  \item the folder you last opened a file from.
\end{itemize}
These are stored in:
\begin{description}[style=nextline, leftmargin=2.2em, font=\sffamily\bfseries]
  \item[macOS] \file{\textasciitilde/Library/Preferences/org.xpreview.XPreview.plist}
  \item[Windows (portable)] \file{settings\textbackslash XPreview\textbackslash XPreview.ini} inside the
        \file{XPreview+} folder
  \item[Linux] \file{\textasciitilde/.config/XPreview/XPreview.conf}
\end{description}
To reset \app{}, quit it and delete that file. On macOS you can instead run
\texttt{defaults delete org.xpreview.XPreview} in Terminal.

\screenshot{window-dracula.png}{\linewidth}{The Modern look with the Dracula dark style.}

\subsection{Privacy and safety}
\begin{itemize}
  \item \app{} never connects to the internet by itself. It contains no analytics, no update
        checks and no accounts.
  \item Documents cannot make \app{} load anything from the internet or from elsewhere on
        your computer. For example, pictures referenced by an SVG or ODT file but stored
        outside it are not loaded.
  \item Macros and scripts in documents are never run, including JavaScript in PDF forms.
  \item PostScript and EPS files are converted by Ghostscript in a sandbox with no network,
        no access to your files, and a time limit.
  \item Passwords you type are not stored or remembered.
  \item Links only open when you click them, and only web and e-mail links.
  \item Every file is treated as potentially hostile. Damaged or deliberately malformed files
        are handled safely and size limits stop them from exhausting memory. Such
        a file shows an error message instead of crashing the application.
  \item Exporting never changes the original file, and edits to a PDF only reach its file
        when you save.
\end{itemize}

% ================================================================
\section{Troubleshooting}

\begin{description}[style=nextline, leftmargin=2.2em, font=\sffamily\bfseries]
  \item[macOS says \app{} cannot be opened or verified]
        Allow it once in \menu{System Settings > Privacy \& Security} (see Section~2.1).
  \item[Windows shows ``Windows protected your PC'']
        Click \ui{More info}, then \ui{Run anyway}.
  \item[\app{} is missing from ``Open With'']
        macOS adds \app{} to that list after it has been copied to \ui{Applications} and opened once.
  \item[A document looks different from the original]
        For RTF and ODT files, click \ui{Details\ldots} in the yellow bar to see what was
        simplified. For an exact rendering, use \ui{Open in Default App}.
  \item[Text cannot be found or selected]
        The document may be a scanned image without a text layer, or an image file.
  \item[``The file is too large to open'']
        \app{} limits very large or deeply nested files to stay safe and responsive. For
        example, images may be at most 30\,000 pixels wide or high.
  \item[A spreadsheet or presentation does not open]
        OpenDocument spreadsheets and presentations are not supported yet.
  \item[HEIC photos do not open on Windows]
        HEIC and JPEG 2000 images are only supported on macOS.
  \item[Something else is wrong]
        Reset \app{} as described in Section~\ref{sec:settings}. If the problem remains, please
        report it (below).
\end{description}

\subsection{Reporting problems}
Please report bugs and ask for features on the project's
\href{https://github.com/swijeyes/XPreview-Plus/issues}{GitHub issues page}. Include your
operating system, the \app{} version (\menu{Help > About \app}), and, if possible, a file that
shows the problem. Report \textbf{security} problems privately through the repository's
\ui{Security} tab instead.

% ================================================================
\section{License and credits}

\app{} is Copyright \textcopyright{} 2026 Sanjeeva Wijeyesakere.

\app{} is free software: you can redistribute it and/or modify it under the terms of the GNU
Affero General Public License as published by the Free Software Foundation, either version 3 of
the License, or (at your option) any later version. It is distributed in the hope that it will be
useful, but \textbf{without any warranty}; without even the implied warranty of merchantability or
fitness for a particular purpose. See the license for details. The full text is included with the
application (\file{LICENSE}). The source code is available at
\url{https://github.com/swijeyes/XPreview-Plus}.

\app{} is built with:
\begin{itemize}
  \item \textbf{Qt} (LGPL-3.0), the application framework;
  \item \textbf{PDFium} (BSD-3-Clause / Apache-2.0), for PDF rendering;
  \item \textbf{miniz} (MIT), for reading OpenDocument packages;
  \item the \textbf{Arimo}, \textbf{Tinos} and \textbf{Cousine} fonts (Apache-2.0), and the
        \textbf{Carlito} and \textbf{Caladea} fonts (SIL Open Font License 1.1).
\end{itemize}
Their licenses are listed in \file{THIRD\_PARTY\_NOTICES.md} and included with the application.
The \app{} icon was created by Sanjeeva Wijeyesakere with Google Gemini.

\end{document}
